\section{How far away must the auxiliary scale be?}
For two error functions $E,F$ the repository selects a point $u$ in $[t,t^2]$ with
$E(u)-X F(u)\le e^{-t}$, uniformly for $0\le X\le e^t$.
Here $u$ is a logarithmic variable. Two earlier ProveIt reports, \emph{Polynomial auxiliary scales and Kummer prime races beyond Dirichlet GRH} (the Kummer window transfer) and the cyclic generator transfer, already obtain linear windows with family-dependent constants and a common averaging measure. We recall their argument using the repository's moment estimate. The supplementary result below gives a necessary and sufficient condition on a schedule that is to work for every polar pair.

\paragraph{The analytic input.}
In \texttt{Lib/Analytic/Landau/PolarLandau.lean}, a real measurable function $E$ has the property \texttt{IsPolar} relative to a residue function if $|E(v)|\le C e^v$ for $v\ge0$ and the integral
\[
  \Phi(s)=\int_0^\infty E(v)e^{(1-s)v}\,dv,\qquad \Re s>2,
\]
has a meromorphic continuation with the specified simple-pole residue limits for $\Re s>0$ and with $(s-1)\Phi(s)$ holomorphic in a neighborhood of every point in $\Re s\ge1$. In this last clause the assigned value of the product at $1$ is zero, so the continuation is removable there.
For any fixed finite list $E_1,\dots,E_m$ of polar functions there are $R>3$ and $B\ge1$ such that for every $j$, $v\ge0$ and integer $k\ge0$,
\begin{equation}
 |E_j(v)|\le B e^v,\qquad
 \left|\int_0^\infty v^k e^{-3v} E_j(v)\,dv\right|
 \le B k! R^{-k}. \label{src:moment}
\end{equation}
This is \texttt{IsPolar.eventually\_moment\_le} in
\texttt{New/Endgame/OneSidedCriterion.lean}. The signed integral is the $k$th Taylor derivative at $s=4$, up to sign. Holomorphy on a neighborhood of the closed disk of radius $3$ about $4$ allows a slightly larger disk for each function. A finite list admits one common radius. The constants may depend on all the functions.

\paragraph{Linear window theorem.}
Assume \eqref{src:moment}, fix $a,b\ge0$ with $a+b>0$, and set $d_0=a+b+2$. Choose
\[
 K>\max\left\{\frac{d_0}{\log(R/3)},\frac{4+d_0}{\log 2}\right\},
 \qquad L>\max\{1,d_0+K\log 3\}.
\]
For every sufficiently large $t$ there is an explicit probability measure $\mu_t$ supported on $[t,Lt]$ for which
\[
 \max_j\left|\int E_j\,d\mu_t\right|\le e^{-(a+b)t}.
\]
Consequently, if $\sum_j |c_j|\le e^{bt}$, some $u\in[t,Lt]$ has
$\sum_j c_j E_j(u)\le e^{-at}$. For continuous $E_j$ one can take
$|\sum_j c_j E_j(u)|\le e^{-at}$. In either pointwise assertion $u$ may depend on the coefficients. For the repository's coefficients $(1,-X)$ with $X\le e^{bt}$ the same conclusion holds: the proof below has sufficient exponential slack to absorb the factor $2$.

\paragraph{Proof.}
Put $n=\lceil Kt\rceil$ and let $Y$ have density
$3^n v^{n-1}e^{-3v}/(n-1)!$ on $v>0$. By \eqref{src:moment},
\[
 |\mathbb E E_j(Y)|\le 3B(3/R)^{n-1}.
\]
The elementary identity $\mathbb E e^{zY}=(3/(3-z))^n$ for real $z<3$ gives the two tail estimates
\[
 \begin{split}
 \mathbb E[e^Y\mathbf1_{Y<t}]
 &\le e^t\mathbb P(Y<t)\le e^{4t}\mathbb E e^{-3Y}
 =e^{4t}2^{-n},\\
 \mathbb E[e^Y\mathbf1_{Y>Lt}]
 &\le e^{-Lt}\mathbb E e^{2Y}=e^{-Lt}3^n .
 \end{split}
\]
Each bound, and the signed full expectation, is $O(e^{-d_0t})$ by the strict choices of $K,L$. As $e^Y\ge1$, the probability that $t\le Y\le Lt$ tends to $1$. Let $\mu_t$ be the law of $Y$ conditional on that event. Subtracting the two tails and dividing by the probability gives $|\int E_j\,d\mu_t|=O(e^{-d_0t})$ uniformly in $j$. A scalar combination with coefficient norm at most $2e^{bt}$ therefore has a mean whose absolute value is less than $e^{-at}$ for all sufficiently large $t$. Some point is below this upper bound. For a continuous scalar combination its image on the interval contains its mean, proving the two-sided assertion.

\paragraph{Sharp universal schedule.}
Let $H(t)\ge t$ eventually. The following are equivalent for fixed $a,b\ge0$ with $a+b>0$.
\begin{enumerate}
\item For every polar pair $E,F$, all sufficiently large $t$, and every $0\le X\le e^{bt}$, some $u\in[t,H(t)]$ satisfies $E(u)-XF(u)\le e^{-at}$.
\item $H(t)/t\to\infty$.
\end{enumerate}
The theorem proves sufficiency since the constant $L$ is fixed for each pair. For necessity suppose $H(t_i)\le L_0t_i$ on a sequence tending to infinity. Choose $0<\delta<\min\{1,(a+b)/L_0\}$. Take $E=0$, $F(v)=-e^{-\delta v}$, $X=e^{bt_i}$. The transform of $F$ is $-1/(s-1+\delta)$, with one simple pole at $1-\delta$, so it satisfies the stated polar conditions. For every $u$ in the proposed interval, $E(u)-XF(u)\ge e^{(b-\delta L_0)t_i}> e^{-at_i}$. This proves necessity even with continuous functions. The counterexample concerns the full analytic class. Arithmetic subfamilies with stronger pole-free regions can have uniform linear constants.

\paragraph{The arithmetic specialization.}
For fixed positive modulus $M$ and a prime $d\equiv7\pmod 8$ with $\gcd(M,d)=1$, \texttt{isPolar\_primeError} in \texttt{New/Endgame/PrimeError.lean} supplies admissible inputs $E(v)=R_{d,M}(e^v)$ and $F(v)=R_{d,1}(e^v)$. Explicitly,
\[
 R_{d,M}(y)=
 \sum_{\substack{n\le y\\n\equiv1\pmod M}}\!
 \Lambda(n)\left(1-\left(\frac nd\right)\right)
 \left(\frac1n-\frac1y\right)-\frac{\log y}{\varphi(M)}-\kappa_{d,M},
\]
where $\kappa_{d,M}=\Re\,\texttt{logDerivMean}(d,M,1)$ in the repository.
The factor $1/n-1/y$ vanishes when $y=n$, giving continuity of each error for $y>0$. In \texttt{eventually\_primeError\_le\_rpow} the parameters satisfy $18480\mid M$ and $d\equiv M/2-1\pmod M$. The function called $X$ above is a nonnegative ratio of two cardinalities, and the code bounds it by $e^t$ with $t=10^6\log x$. With $a=b=1$, the auxiliary scale from the prior reports $x_1=e^u$ satisfies
\[
  x^{10^6}\le x_1\le x^{10^6 L},
  \qquad |R_{d,M}(x_1)-X R_{d,1}(x_1)|\le x^{-10^6}.
\]
The previous choice allowed an upper endpoint $\exp((10^6\log x)^2)$. For fixed $d,M$, $L$ is independent of $x$ and $X$. This strengthens the internal scale selection. The later determinant and rank estimates still govern the final exponent of the prime error.
